% !TeX root = ../../elegantbook-cn.tex


\section{函数单调性}
函数的单调性是函数最重要的性质。面对较为复杂的函数时，通常需要借助导数研究其单调性。

对于高中阶段常见的可导函数$f(x)$，若在某一区间内恒有$f'(x)>0$，则$f(x)$在该区间内单调递增；若恒有$f'(x)<0$，则$f(x)$在该区间内单调递减。

若$f'(x)$在$x_0$左侧为正、右侧为负，即$f(x)$在$x_0$两侧先增后减，则$x_0$为$f(x)$的极大值点；若$f'(x)$在$x_0$左侧为负、右侧为正，即$f(x)$在$x_0$两侧先减后增，则$x_0$为$f(x)$的极小值点。

\begin{remark}
只有$f'(x_0)=0$并不能判断$x_0$是否为极值点，还需要考察$f'(x)$在$x_0$两侧的符号变化。
\end{remark}

\subsection{不含参函数的单调性}

不含参数的函数单调性问题，核心任务是判断导函数$f'(x)$的正负。题目虽然研究的是$f(x)$的增减变化，但真正需要处理的对象通常是$f'(x)$。因此，解决这类问题时，要先完成下面的转化：
\[
\boxed{
\text{研究$f(x)$的单调性}
\quad\Longrightarrow\quad
\text{研究$f'(x)$的符号}
}
\]

\medskip
\noindent{\bfseries\color{structurecolor}一、基本步骤}

求一个不含参数的函数的单调区间，一般按照以下步骤进行。

\begin{enumerate}
  \item \textbf{求定义域。} 单调区间必须是函数定义域的子区间，因此求导之前应先明确函数的定义域。
  \item \textbf{求导函数。} 计算$f'(x)$，并尽量将其因式分解或整理成便于判断符号的形式。
  \item \textbf{优先整体判号。} 先观察在整个定义域内是否恒有$f'(x)\geqslant0$或$f'(x)\leqslant0$。若能够直接确定导函数恒为非负或恒为非正，就不必再划分区间，可以直接判断函数在相应区间内单调递增或单调递减。若要说明函数严格单调，还需保证$f'(x)$不在任何一个子区间内恒等于$0$。
  \item \textbf{寻找分界点。} 求出定义域内使$f'(x)=0$的点，同时注意使$f'(x)$不存在但$f(x)$有定义的点。
  \item \textbf{判断导函数的符号。} 用分界点将定义域划分为若干区间，分别判断每个区间内$f'(x)$的正负。
  \item \textbf{写出单调区间。} 在$f'(x)>0$的区间上，$f(x)$单调递增；在$f'(x)<0$的区间上，$f(x)$单调递减。
\end{enumerate}

整个过程可以概括为
\[
\boxed{
\text{定义域}
\longrightarrow
\text{求导}
\longrightarrow
\text{优先整体判号}
\longrightarrow
\text{找分界点}
\longrightarrow
\text{判断$f'(x)$的符号}
\longrightarrow
\text{写出单调区间}
}
\]

\medskip
\noindent{\bfseries\color{structurecolor}二、导函数的符号判断}

求导通常不是这类问题最困难的部分，真正的关键是判断$f'(x)$的符号。得到导函数后，不要急于直接讨论，应先观察它的结构，再选择合适的处理方法。

\medskip
\noindent{\bfseries\color{structurecolor}1. 因式分解}

若$f'(x)$是多项式，或能够整理为若干因式的乘积，应优先进行因式分解。例如
\[
f'(x)=(x-a)(x-b)g(x),
\]
其中$g(x)$的符号已经确定，那么$f'(x)$的符号就由$x-a$与$x-b$决定。此时以$a,b$为分界点，利用各因式的符号即可判断$f'(x)$的正负。

对于分式形式
\[
f'(x)=\frac{g(x)}{h(x)},
\]
不仅要考察分子为$0$的点，还要结合定义域排除分母为$0$的点，再分别判断分子、分母的符号。

\medskip
\noindent{\bfseries\color{structurecolor}2. 提取恒正因式}

若导函数中含有指数函数、平方项或其他恒正部分，可以先将这些因式提取出来。例如
\[
f'(x)=\mathrm{e}^x g(x),
\]
由于$\mathrm{e}^x>0$恒成立，因此$f'(x)$与$g(x)$同号。原来的符号判断便转化为判断$g(x)$的符号。

处理这类式子时，应始终关注：\textbf{哪些部分的符号已经确定，哪些部分才真正决定整个导函数的符号。}

\medskip
\noindent{\bfseries\color{structurecolor}3. 再构造函数}

有时$f'(x)=0$无法直接求解，或者$f'(x)$的符号难以直接判断。此时可以把导函数中较复杂的部分设为一个新函数，例如令
\[
g(x)=f'(x),
\]
再通过研究$g'(x)$的符号，判断$g(x)$的单调性，进而确定$g(x)$的零点个数及正负情况。

这一过程实际上是连续的两次转化：
\[
\boxed{
\text{$f(x)$的单调性}
\longrightarrow
\text{$f'(x)$的符号}
\longrightarrow
\text{$f'(x)$的单调性}
}
\]

需要注意的是，研究$g(x)=f'(x)$的单调性只是手段，最终仍要回到$g(x)=f'(x)$的正负，才能得到原函数$f(x)$的单调区间。

\medskip
\noindent{\bfseries\color{structurecolor}三、单调性表的书写}

找到所有分界点后，可以借助表格整理导函数的符号与原函数的变化情况。若$x_1<x_2$为两个分界点，单调性表可以写成如下形式：
\[
\begin{array}{c|ccccccc}
x      & (-\infty,x_1) & x_1 & (x_1,x_2) & x_2 & (x_2,+\infty) \\ \hline
f'(x)  & +              & 0   & -           & 0   & +               \\ \hline
f(x)   & \nearrow       &     & \searrow    &     & \nearrow
\end{array}
\]

由表可知，$f(x)$在$(-\infty,x_1)$和$(x_2,+\infty)$上单调递增，在$(x_1,x_2)$上单调递减；同时，$x_1$是极大值点，$x_2$是极小值点。

\begin{remark}
单调区间通常写成定义域内的区间，不能跨过函数无定义的点。若两个递增区间被一个无定义点隔开，也不能将它们合并为一个区间。
\end{remark}


因此，不含参数的函数单调性问题虽然形式多样，但主线始终不变：先把原函数的增减问题转化为导函数的符号问题，再通过因式分解、提取恒正因式或构造新函数等方法完成符号判断。


\subsection{含参函数的单调性}

含参函数的单调性问题与不含参问题的主线完全相同，仍然是通过判断$f'(x)$的符号研究$f(x)$的单调性。二者的区别在于：参数的取值会改变导函数零点的个数、位置和大小关系，因此必须根据参数进行分类讨论。

这类问题的核心不是“分类越多越好”，而是找到真正会使单调性发生变化的参数临界值。整体思路可以概括为
\[
\boxed{
\text{求导并化简}
\longrightarrow
\text{寻找决定符号的部分}
\longrightarrow
\text{确定参数临界值}
\longrightarrow
\text{分类判断$f'(x)$的符号}
\longrightarrow
\text{写出单调区间}
}
\]

\medskip
\noindent{\bfseries\color{structurecolor}一、含参问题需要讨论什么？}

含参问题中的分类讨论，通常来自以下四个方面。

\begin{enumerate}
  \item \textbf{最高次项系数是否为$0$。} 若导函数中某个含参数的系数可能为$0$，则参数取到该值时，导函数的次数或结构会发生变化，必须单独讨论。
  \item \textbf{导函数零点的个数。} 当$f'(x)=0$可转化为一元二次方程时，判别式$\Delta$的正负决定零点有两个、一个还是没有，因此$\Delta=0$往往是重要的参数临界情况。
  \item \textbf{零点是否在定义域内。} 即使已经求出$f'(x)=0$的根，也要判断这些根是否属于函数的定义域。根越过定义域端点时，单调区间也会发生变化。
  \item \textbf{多个分界点的大小关系。} 如果导函数的零点为$x_1(a),x_2(a)$，或者还存在固定的分界点$x_0$，就需要比较它们的大小。两个分界点重合或交换顺序时，对应的参数值就是分类标准。
\end{enumerate}

因此，参数临界值通常来自
\[
\boxed{
\text{系数为$0$、判别式为$0$、根等于定义域端点、两个分界点相等}
}
\]

\medskip
\noindent{\bfseries\color{structurecolor}二、基本处理步骤}

\begin{enumerate}
  \item \textbf{确定定义域。} 若定义域本身含有参数，应先讨论定义域是否为空以及端点的大小关系。
  \item \textbf{求导并化简。} 计算$f'(x)$，尽量进行因式分解，或提取恒正、恒负的因式，找出真正决定符号的部分。
  \item \textbf{优先整体判号。} 先判断是否存在某些参数范围，使$f'(x)\geqslant0$或$f'(x)\leqslant0$在整个定义域内恒成立。若可以整体判号，就不必继续求零点和划分区间。
  \item \textbf{求导函数的零点。} 在不能整体判号时，解$f'(x)=0$，并检查所得根是否属于定义域。
  \item \textbf{确定分类标准。} 根据零点个数、零点与定义域端点的位置关系以及各分界点的大小关系，找出参数临界值。
  \item \textbf{分类列出单调区间。} 在每一种参数范围内固定分界点的顺序，判断$f'(x)$的符号，最后写出$f(x)$的单调递增区间和单调递减区间。
\end{enumerate}

\begin{remark}
分类讨论应遵循“不重不漏”的原则。各类参数范围的并集应覆盖参数的全部允许取值，不同类别之间不能重复。
\end{remark}

\medskip
\noindent{\bfseries\color{structurecolor}三、常见的导函数结构}

\medskip
\noindent{\bfseries\color{structurecolor}1. 已经因式分解的结构}

若导函数可以整理为
\[
f'(x)=(x-x_0)\bigl(x-\varphi(a)\bigr)g(x),
\]
且$g(x)$的符号已经确定，那么单调性主要由两个分界点$x_0$与$\varphi(a)$决定。此时不必进行没有依据的分类，只需比较
\[
\varphi(a)<x_0,qquad
\varphi(a)=x_0,qquad
\varphi(a)>x_0.
\]
也就是说，分类标准来自方程$\varphi(a)=x_0$。在每一种情况下确定两个分界点的先后顺序后，再画符号图即可。

如果某个因式的次数为偶数，例如$(x-x_0)^2$，则该因式经过$x_0$时不变号；如果次数为奇数，则经过$x_0$时变号。判断导函数符号时，不仅要找零点，还要注意零点的重数。

\medskip
\noindent{\bfseries\color{structurecolor}2. 关于$x$的一元二次式}

若导函数的符号由
\[
g(x)=A(a)x^2+B(a)x+C(a)
\]
决定，应先讨论$A(a)=0$的情况，因为此时$g(x)$不再是二次函数。对于$A(a)\neq0$的情况，再计算判别式
\[
\Delta=B^2(a)-4A(a)C(a).
\]

讨论顺序通常为
\[
\boxed{
A(a)=0
\quad\text{或}\quad
\begin{cases}
A(a)\neq0,\ \Delta<0,\\
A(a)\neq0,\ \Delta=0,\\
A(a)\neq0,\ \Delta>0.
\end{cases}
}
\]

当$\Delta<0$时，$g(x)$的符号由二次项系数$A(a)$决定，往往可以直接得到原函数在整个定义域内的单调性；当$\Delta=0$时，导函数只有一个二重零点，经过该点时通常不变号；当$\Delta>0$时，导函数有两个不同的零点，需要结合开口方向判断各区间内的符号。

\medskip
\noindent{\bfseries\color{structurecolor}3. 参数与自变量可以分离的结构}

若$f'(x)=0$难以直接解出$x$，但可以整理为
\[
a=\varphi(x),
\]
则可以先研究函数$y=\varphi(x)$的图像、单调性和值域。方程$f'(x)=0$的根的个数，就转化为水平直线$y=a$与曲线$y=\varphi(x)$的交点个数。

此时，$\varphi(x)$的极值就是参数$a$的临界值。参数跨过这些临界值时，交点个数发生变化，导函数的零点个数也随之变化。这种方法本质上是将“含参方程的根”转化为“函数图像的交点”。

类似地，若导函数能够整理为
\[
f'(x)=g(x)-a,
\]
则判断$f'(x)$的正负等价于比较$g(x)$与$a$的大小：
\[
f'(x)>0\Longleftrightarrow g(x)>a,
\qquad
f'(x)<0\Longleftrightarrow g(x)<a.
\]
因此，可以先研究$g(x)$的单调性与值域，再利用直线$y=a$进行分类。

\medskip
\noindent{\bfseries\color{structurecolor}四、已知单调性求参数范围}

有些题目不是要求讨论函数的单调区间，而是已知$f(x)$在某个区间$I$上单调，要求参数的取值范围。此时应将单调性翻译为导函数在区间上的恒成立问题：
\[
\begin{aligned}
f(x)\text{在$I$上单调递增}
&\Longrightarrow f'(x)\geqslant0\quad(x\in I),\\
f(x)\text{在$I$上单调递减}
&\Longrightarrow f'(x)\leqslant0\quad(x\in I).
\end{aligned}
\]

这样，原问题就转化为含参不等式在区间$I$上恒成立。处理时可以根据导函数的结构，选择以下思路：

\begin{enumerate}
  \item 分离参数，将不等式整理为$a\geqslant\varphi(x)$或$a\leqslant\varphi(x)$，再求$\varphi(x)$在区间$I$上的最值；
  \item 利用二次函数在给定区间上的最值，保证$f'(x)$恒为非负或恒为非正；
  \item 根据导函数零点的位置，要求零点落在区间外部或恰好位于区间端点。
\end{enumerate}

需要注意的是，由“函数单调”转化为“导函数恒为非负或恒为非正”时，端点和等号往往可以取到，不能机械地写成$f'(x)>0$或$f'(x)<0$，否则可能漏掉临界参数。

\medskip
\noindent{\bfseries\color{structurecolor}五、分类讨论的书写原则}

含参函数单调性问题最容易出现的错误，是还没有找到分类标准就开始大量讨论。规范的做法是：先明确参数改变了什么，再根据发生变化的位置分类。

书写时可以采用以下顺序：
\[
\boxed{
\text{先给出参数临界值}
\longrightarrow
\text{再说明每类中分界点的顺序}
\longrightarrow
\text{最后写出对应的单调区间}
}
\]

每讨论完一种参数范围，都应检查以下三个问题：定义域是否正确，分界点是否都在定义域内，各分界点的顺序是否已经确定。只要这三点明确，导函数的符号和原函数的单调区间就能顺利写出。
